% !TeX root = ../教材题库_全部条目_转写工作稿.tex
% Source: XJ-B1-C06-SU-P262-Y1246
% 第6章 / 小结与复习 小结与复习 / 教材 p.262 / PDF 第269页
% OCR未校对：公式与文字均可能有误；未完成数学TeX转写；禁止直接发给学生
% 原题及相邻原解截图（核对用）：../题目原图/XJ-B1-C06-SU-P262-Y1246-01.webp
% 本文件只在转写完成后作为题干使用；OCR 条目尚不是可靠题干。
\begin{bankitem}{第6章 小结与复习 第2题}{第6章 / 小结与复习 小结与复习 / 教材 p.262 / PDF 第269页}{XJ-B1-C06-SU-P262-Y1246}{OCR未校对：公式与文字均可能有误；未完成数学TeX转写；禁止直接发给学生}
\begin{ocrdraft}
[教材p.262 / PDF第269页]
\par
2．如果我们能够掌握总体的全部数据，那么只需做一些简单的统计描述，就
\par
可得到所关心的总体特征，如总体均值、方差等．但现实情况比较复杂，有些现
\par
象的范围比较广，不可能对总体中的每个个体进行测量，或者总体的个体数据多，
\par
不可能也没有必要进行一一测量，这就需要从总体中随机抽取一部分个体进行调
\par
查，抽样调查是进行统计调查的主要方式．在抽取样本的过程中，要遵循的最基
\par
本的原则是什么？简单随机抽样与分层抽样各有何特点？
\end{ocrdraft}
\end{bankitem}
